% Part:sets-functions-relations
% Chapter: size-of-sets
% Section: enumerations-alt

\documentclass[../../../include/open-logic-section]{subfiles}

\begin{document}

\olfileid{sfr}{siz}{enm-alt}

\olsection{گنتی وار فہرستاں تے \usetoken{S}{enumerable} سیٹ}

\begin{editorial}
  ایہ حصہ گنتی وار فہرستاں نوں $\Nat$ (دے مُڈھلے ٹکڑیاں) نال
  اک بہ اک مطابقتاں دے طور تے تعریف کردا اے، جیویں سیٹ تھیوری وچ
  کیتا جاندا اے۔ ایس لئی ایہ \olref[enm]{sec} دیاں تعریفاں نال ذرا
  ٹکراؤندا اے، تے اوتھے دیاں ساریاں مثالاں دُہراؤندا اے۔ ایہ اوس
  حصے نالوں کجھ مختصر وی اے۔
\end{editorial}

اسی کسے متناہی سیٹ دے !!{element}s نوں بس ترتیب وار لکھ کے اوہ سیٹ
دس سکدے آں۔ اسی ایویں کردے آں جدوں کوئی سیٹ ایس طرح تعریف کریے:
\[
  A = \{a_1, a_2, \ldots, a_n\}.
\]
% PNBSIZ-003 ماخذ دی درستی: انگریزی ماخذ دا پہلا جملہ نحوی طور تے ٹُٹیا ہویا اے؛ اگلی نمائش تے وضاحت موجب مراد متناہی سیٹ دے رکناں نوں ترتیب وار لکھ کے اوہنوں دسنا اے۔
جے فرض کریے کہ !!{element}s $a_1$، \dots، $a_n$ سارے وکھرے نیں،
تاں ایہ ساڈے لئی~$A$ تے پہلے $n$ قدرتی عدداں $0$، \dots، $n-1$
دے وچکار !!a{bijection} دیندا اے۔ الٹ پاسے، کیوں جے ہر متناہی سیٹ
وچ صرف متناہی گنتی دے !!{element}s ہوندے نیں، ہر متناہی سیٹ نوں
ایہو جہی مطابقت وچ رکھیا جا سکدا اے۔ ہور لفظاں وچ، جے~$A$ متناہی
اے، تاں~$A$ تے $\{0, \dots, n-1\}$ دے وچکار !!a{bijection} موجود
اے، جتھے~$n$،~$A$ دے !!{element}s دی گنتی اے۔

جے اسی کجھ قسماں دے لامتناہی سیٹاں نوں وی آون دیواں، تاں اسی کجھ
لامتناہی سیٹاں دی وی گنتی وار فہرست بنن دی اجازت دیواں گے۔ اسی ایہ
گل ایویں ٹھیک طور تے کہہ سکدے آں کہ اک لامتناہی سیٹ دی گنتی وار
فہرست اوس سیٹ تے سارے~$\Nat$ دے وچکار !!a{bijection} بناؤندی اے۔

\begin{defn}[گنتی وار فہرست، سیٹ تھیوری دے لحاظ نال]
سیٹ~$A$ دی اک \emph{گنتی وار فہرست} اوہ !!a{bijection} اے جیہدی
حاصل قدراں دا سیٹ~$A$ ہووے تے جیہدا دائرۂ تعریف یا قدرتی عدداں دا
اک مُڈھلا سیٹ $\{0, 1, \ldots, n\}$ {یا} قدرتی عدداں دا پورا
سیٹ~$\Nat$ ہووے۔
\end{defn}

\begin{explain}
\emph{گنتی وار فہرست} لفظ دے ایس ورتاوے تھلے اک بدیہی خیال اے۔ ایہ
آکھنا کہ اسی سیٹ~$A$ دی گنتی وار فہرست بنا لئی اے، ایہ آکھنا اے کہ
کوئی !!a{bijection}~$f$ موجود اے جیہڑا سانوں سیٹ~$A$ دے عنصراں نوں
گن کے کڈھن دیندا اے۔ $0$واں عنصر $f(0)$ اے، پہلا $f(1)$ اے،
\ldots{}، $n$واں $f(n)$ اے، \ldots۔\footnote{ہاں، اسی گنتی
$0$ توں شروع کردے آں۔ بے شک اسی~$1$ توں وی شروع کر سکدے ساں۔ ایس
نال کوئی وڈا فرق نہ پیندا۔ سانوں بس~$\Nat$ دی تھاں~$\PosInt$ رکھنا
پیندا۔} ایس دی وجہ اگلی گل نوں شامل کر کے ہور وی صاف کیتی جا سکدی اے:
\end{explain}

\begin{defn}
  \ollabel{defn:enumerable}
  سیٹ~$A$، !!{enumerable} عین اوس ویلے اے جدوں یا تاں
  $A = \emptyset$ ہووے یا~$A$ دی کوئی گنتی وار فہرست موجود ہووے۔
  اسی~$A$ نوں !!{nonenumerable} عین اوس ویلے آکھدے آں جدوں~$A$،
  !!{enumerable} نہ ہووے۔
\end{defn}

\begin{explain}
سو اک سیٹ !!{enumerable} عین اوس ویلے اے جدوں اوہ خالی ہووے یا تسیں
اک گنتی وار فہرست نال اوہدے !!{element}s گن کے کڈھ سکو۔
\end{explain}

\begin{ex}
قدرتی عدداں دی گنتی وار فہرست بناؤن والا فنکشن بس شناختی فنکشن
$\Id{\Nat} \colon \Nat \to \Nat$ اے، جیہڑا $\Id{\Nat}(n) = n$ نال
دتا گیا اے۔ \emph{مثبت} قدرتی عدداں، $\Nat^+ =
\Nat \setminus \{0\}$، دی گنتی وار فہرست بناؤن والا فنکشن
$g(n) = n + 1$ اے، یعنی جانشین فنکشن۔
\end{ex}

\begin{prob}
وکھاؤ کہ سیٹ~$A$، !!{enumerable} عین اوس ویلے اے جدوں یا تاں
$A = \emptyset$ ہووے یا کوئی !!a{surjection}
$f\colon \Nat \to A$ موجود ہووے۔ وکھاؤ کہ~$A$، !!{enumerable}
عین اوس ویلے اے جدوں کوئی !!a{injection} $g\colon A \to \Nat$
موجود ہووے۔
\end{prob}

\begin{ex}
فنکشن $f\colon \Nat \to \Nat$ تے $g \colon \Nat \to \Nat$، جیہڑے
ایویں دتے گئے نیں:
\begin{align*}
  f(n) & = 2n \text{ تے}\\
  g(n) & = 2n+1
\end{align*}
ترتیب وار جفت قدرتی عدداں تے طاق قدرتی عدداں دیاں گنتی وار فہرستاں
بناؤندے نیں۔ پر دوہاں وچوں کوئی وی !!{surjective} نہیں، سو کوئی وی
$\Nat$ دی گنتی وار فہرست نہیں۔
\end{ex}

\begin{prob}
مربع عدداں $1$، $4$، $9$، $16$، \dots دی گنتی وار فہرست تعریف کرو۔
\end{prob}

\begin{ex}
$\lceil x \rceil$ نوں \emph{چھت} فنکشن منّو، جیہڑا $x$ نوں اُتے ول
سب توں نیڑے صحیح عدد تک پورا کردا اے۔ فیر فنکشن
$f \colon \Nat \to \Int$، جیہڑا ایویں دتا گیا اے:
\[
  f(n) = (-1)^{n} \left\lceil\tfrac{n}{2}\right\rceil
\]
صحیح عدداں دے سیٹ~$\Int$ دی گنتی وار فہرست اگلے ڈھنگ نال بناؤندا اے:
\[
\begin{array}{c c c c c c c c}
f(0) & f(1) & f(2) & f(3) & f(4) & f(5) & f(6) & \dots \\ \\
\big\lceil \tfrac{0}{2} \big\rceil & -\big\lceil \tfrac{1}{2}\big\rceil &  \big\lceil \tfrac{2}{2} \big\rceil & -\big\lceil \tfrac{3}{2} \big\rceil & \big\lceil \tfrac{4}{2} \big\rceil  & -\big\lceil \tfrac{5}{2}\big\rceil & \big\lceil \tfrac{6}{2} \big\rceil & \dots \\ \\
0 & -1 & 1 & -2 & 2 & -3 & 3& \dots
\end{array}
\]
دھیان دیو کہ $f$، مثبت تے منفی صحیح عدداں دے وچکار ”ٹپدے“ ہوئے
$\Int$ دیاں قدراں پیدا کردا اے۔ تسیں~$f$ نوں اگلی صورتاں نال تعریف
کیتا ہویا وی سمجھ سکدے او:
\[
f(n) = \begin{cases}
  \frac{n}{2} & \text{جے $n$ جفت اے}\\
  -\frac{n+1}{2} & \text{جے $n$ طاق اے}
  \end{cases}
\]
\end{ex}

\begin{prob}
وکھاؤ کہ جے $A$ تے~$B$، !!{enumerable} نیں، تاں $A \cup B$ وی اے۔
\end{prob}

\begin{prob}
$n$ اُتے استقرا نال وکھاؤ کہ جے $A_1$، $A_2$، \dots، $A_n$ سارے
!!{enumerable} نیں، تاں $A_1 \cup \dots \cup A_n$ وی اے۔
\end{prob}

\end{document}
